\documentclass[12pt]{article}
\usepackage[utf8]{inputenc}
\usepackage[T1]{fontenc}
\usepackage{hyperref}
\title{Tlingit Language Resources for K-12 Curriculum Development}
\author{Compiled by Claude Sonnet 5 (claude-sonnet-5) for Andr\'as Kornai}
\date{}
\begin{document}
\maketitle

\section{Introduction}
Tlingit (ISO 639-3 tli) is an endangered language of southeast Alaska and neighbouring Canada. The reference description says it is not known to be taught in schools. The search results show otherwise: Sealaska Heritage Institute (SHI) has produced school curricula, and Juneau has a Tlingit programme.

\section{Primary Educational Resources}
\begin{itemize}
\item SHI Tlingit Language Proficiency Scope and Sequence, written by birth speakers, teachers, learners and administrators, with four levels: \url{https://sealaskaheritage.org/?p=10159}
\item SHI's catalogue of teaching and learning resources: \url{https://sealaskaheritage.org/alaskan-native-language-learners/shi-southeast-alaskan-native-language-learning-teaching-resources/}
\item An 18-unit Tlingit language and culture curriculum made by SHI and the Juneau School District, reported in 2007: \url{https://listserv.linguistlist.org/pipermail/ilat/2007-July/003257.html} and \url{https://listserv.linguistlist.org/pipermail/ilat/2007-July/003290.html}
\end{itemize}

\section{Linguistic and Reference Materials}
The Alaska State Library page describes SHI's books, workbooks, apps, games, audio and video: \url{https://lam.alaska.gov/sled/languages/aknative}. A Tlingit verb database is linked from \url{https://sfthcc.org/learn-tlingit-haida-or-tsimshian}.

\section{Religious and Community Resources}
SHI language department: \url{https://sealaskaheritage.org/shi-language-department/}. The Tlingit Culture, Language and Literacy programme in Juneau began with SHI in 2000: \url{https://sealaskaheritage.org/?p=5383}. Goldbelt Heritage Foundation resources are linked from the sfthcc.org page above.

\section{Digital and Audio Resources}
SHI offers the Learning Tlingit app (vocabulary with audio) and a Language Games app for Tlingit, Haida and Tsimshian: \url{https://sealaskaheritage.org/?p=10159}.

\section{Curriculum Development Context}
Tlingit has the most developed school resources of the four languages in this batch. The search results did not show the current availability of the 2007 curriculum, so check SHI directly.
\end{document}